\documentclass[a4paper, serif, 10pt]{moderncv}

%% moderncv
% color and style
\moderncvcolor{black}
\moderncvstyle{banking}

%% page layout/margins
\usepackage[top=2.5cm, bottom=2.5cm, left=1.5cm, right=1.5cm]{geometry}

\nopagenumbers{}

%% Babel config for French language
% ref:
% - https://latex3.github.io/babel/guides/locale-french.html
\usepackage[french]{babel}

%% fontspec
\usepackage{fontspec}

\IfFontExistsTF{Linux Libertine}{
  \setmainfont[Ligatures={TeX, Common}, Numbers=OldStyle]{Linux Libertine}
}{}
\IfFontExistsTF{Linux Libertine O}{
  \setmainfont[Ligatures={TeX, Common}, Numbers=OldStyle]{Linux Libertine O}
}{}

%% CTeX
% CJK support, used to show CJK characters in the resume
%
% - heading=false: disable ctex heading
% - punct=kaiming: use kaiming punctuations styles for CJK
% - scheme=plain: use plain scheme, do not override \normalsize font size
% - space=auto: space settings for CJK characters
%
% CJK fonts are resolved per locale language, e.g. Noto Serif/Sans CJK JP
% for Japanese resumes, with fallback to the next region if not installed
%
% ref:
% - http://ctan.mirrorcatalogs.com/language/chinese/ctex/ctex.pdf
\usepackage[UTF8, heading=false, punct=kaiming, scheme=plain, space=auto]{ctex}

\IfFontExistsTF{Noto Serif CJK SC}{
  \setCJKmainfont{Noto Serif CJK SC}
}{}
\IfFontExistsTF{Noto Sans CJK SC}{
  \setCJKsansfont{Noto Sans CJK SC}
}{}

%% line spacing
\usepackage{setspace}
\setstretch{1.125}

%% URL styling - use normal text instead of monospace
\urlstyle{same}

%% Patch \httplink and \httpslink to handle full URLs
% The original moderncv macros always prepend http:// or https:// to URLs.
% This causes issues when the URL already has a protocol (e.g., https://example.com)
% resulting in malformed URLs like https://https://example.com.
% This patch checks if the URL already contains "://" and uses it directly if so.
% Note: We use \str_set:Nx (x-expansion) to expand macros like \@homepage first.
\ExplSyntaxOn
\str_new:N \g_yamlresume_url_str

\RenewDocumentCommand{\httpslink}{O{}m}{
  \str_set:Nx \g_yamlresume_url_str {#2}
  \str_if_in:NnTF \g_yamlresume_url_str {://}
    { \url{#2} }
    { \url{https://#2} }
}

\RenewDocumentCommand{\httplink}{O{}m}{
  \str_set:Nx \g_yamlresume_url_str {#2}
  \str_if_in:NnTF \g_yamlresume_url_str {://}
    { \url{#2} }
    { \url{http://#2} }
}
\ExplSyntaxOff

\name{Camille Moreau}{}
\title{Responsable des opérations}
\phone[mobile]{+33 6 12 34 56 78}
\email{camille.moreau@example.fr}
\homepage{https://www.camille-moreau.fr}

\address{24 rue des Archives, Paris, Île-de-France, France, 75004}{}{}

\extrainfo{{\small \faLinkedin}\ \href{https://www.linkedin.com/in/camille-moreau-ops}{@camille-moreau-ops} {} {} {} • {} {} {} 
{\small \faTwitter}\ \href{https://twitter.com/camille\_ops}{@camille\_ops}}

\begin{document}

\maketitle

\section{Informations de base}

\cvline{}{\begin{itemize}
\item Responsable des opérations avec plus de 10 ans d'expérience dans la logistique, la supply chain et le pilotage de sites.
\item Solide expertise en management d'équipes pluridisciplinaires, en amélioration continue et en optimisation des processus.
\item Capacité démontrée à piloter des budgets, des indicateurs de performance et des projets transverses.
\item Anglais professionnel et pratique courante des outils de reporting et d'analyse de données.
\end{itemize}}
\section{Formation}

\cventry{sept. 2012–juin 2014}
        {Master, Management des opérations et de la supply chain, Score : Mention Bien}
        {Université Paris-Dauphine}
        {\url{https://www.dauphine.psl.eu}}
        {}
        {\begin{itemize}
\item Formation solide en pilotage des opérations et en gestion des flux.
\item Projets académiques consacrés à l'optimisation logistique et à la performance industrielle.
\item Développement de compétences en analyse quantitative et en management d'équipe.
\end{itemize}
\textbf{Cours} : Gestion des opérations, Supply Chain Management, Contrôle de gestion, Gestion de projet, Analyse de données, Management d'équipe}
\section{Expérience professionnelle}

\cventry{janv. 2020–Present}
        {Responsable des opérations}
        {Groupe Horizon}
        {\url{https://www.groupehorizon.fr}}
        {}
        {\begin{itemize}
\item Pilotage quotidien des opérations d'un site logistique de 120 collaborateurs.
\item Management direct de 18 personnes et animation d'une équipe transverse de 6 responsables.
\item Déploiement d'indicateurs de performance et de tableaux de bord hebdomadaires.
\item Réduction des coûts opérationnels de 12 \% en deux ans grâce à la refonte des processus.
\item Coordination des relations avec les transporteurs, fournisseurs et clients grands comptes.
\end{itemize}
\textbf{Mots-clés} : Pilotage opérationnel, Management d'équipe, Amélioration continue, KPI, Budget}

\cventry{mars 2017–déc. 2019}
        {Chef d'équipe logistique}
        {LogiFret Solutions}
        {\url{https://www.logifret-solutions.fr}}
        {}
        {\begin{itemize}
\item Encadrement d'une équipe de 25 préparateurs et caristes sur deux shifts.
\item Mise en place d'un plan de formation et de polyvalence pour réduire les ruptures de service.
\item Suivi des stocks, des expéditions et de la qualité de service client.
\item Participation active à la certification ISO 9001 du site.
\end{itemize}
\textbf{Mots-clés} : Logistique, Planification, Sécurité, Performance, Qualité}

\cventry{sept. 2014–févr. 2017}
        {Coordinateur des opérations}
        {Maison Delcourt}
        {\url{https://www.maison-delcourt.fr}}
        {}
        {\begin{itemize}
\item Coordination des approvisionnements et des livraisons pour trois sites régionaux.
\item Optimisation des tournées de transport et réduction des délais de livraison.
\item Suivi budgétaire des achats et négociation avec les fournisseurs.
\item Amélioration de la satisfaction client par la mise en place d'un suivi des réclamations.
\end{itemize}
\textbf{Mots-clés} : Coordination, Approvisionnement, Stocks, Transport, Relation fournisseurs}
\section{Langues}

\cvline{Français}{Niveau natif ou bilingue}
\cvline{Anglais}{Niveau professionnel complet \hfill \textbf{Mots-clés} : TOEIC 920}
\cvline{Espagnol}{Niveau de travail limité \hfill \textbf{Mots-clés} : Niveau B1}
\section{Compétences}

\cvline{Gestion des opérations}{Expert \hfill \textbf{Mots-clés} : Planification, Organisation, Processus, Qualité}
\cvline{Management d'équipe}{Expert \hfill \textbf{Mots-clés} : Recrutement, Coaching, Animation, Évaluation}
\cvline{Amélioration continue}{Avancé \hfill \textbf{Mots-clés} : Lean, Six Sigma, Kaizen, PDCA}
\cvline{Supply chain}{Avancé \hfill \textbf{Mots-clés} : Approvisionnement, Stocks, Transport, Logistique}
\cvline{Analyse et pilotage}{Intermédiaire \hfill \textbf{Mots-clés} : KPI, Tableaux de bord, Excel, Power BI}
\section{Récompenses}

\cventry{mars 2023}
        {Groupe Horizon}
        {Prix de l'innovation opérationnelle}
        {}
        {}
        {Distinction attribuée pour la refonte d'un processus logistique ayant permis de réduire les coûts et les délais de traitement.}
\section{Certificats}

\cventry{juin 2021}
        {IASSC}
        {Ceinture Verte Lean Six Sigma}
        {\url{https://www.iassc.org}}
        {}
        {}

\cventry{janv. 2022}
        {ASCM}
        {Certification en gestion de la supply chain}
        {\url{https://www.ascm.org}}
        {}
        {}
\section{Projets}

\cventry{janv. 2021–nov. 2021}
        {Migration et déploiement d'un ERP pour les opérations logistiques et achats.}
        {Déploiement ERP}
        {\url{https://www.groupehorizon.fr/projets/erp}}
        {}
        {\begin{itemize}
\item Coordination des équipes métiers, informatiques et fournisseurs.
\item Cartographie et simplification des processus avant migration.
\item Plan de formation des utilisateurs et accompagnement du changement.
\end{itemize}
\textbf{Mots-clés} : ERP, Conduite du changement, Processus, Formation}

\cventry{avr. 2018–oct. 2018}
        {Refonte des tournées de livraison pour améliorer la ponctualité et réduire les coûts.}
        {Optimisation des tournées}
        {\url{https://www.logifret-solutions.fr/projets/tournees}}
        {}
        {\begin{itemize}
\item Analyse des données de transport et des contraintes clients.
\item Nouveau découpage des zones et ajustement des fréquences de livraison.
\item Réduction des kilomètres parcourus et amélioration du taux de service.
\end{itemize}
\textbf{Mots-clés} : Optimisation, Tournées, Coûts, Délais}
\section{Centres d'intérêt}

\cvline{Randonnée}{Sentiers, Nature, Photographie}
\cvline{Cuisine}{Cuisine française, Pâtisserie, Œnologie}
\section{Bénévolat}

\cventry{oct. 2019–Present}
        {Bénévole logistique}
        {Les Restos du Cœur}
        {\url{https://www.restosducoeur.org}}
        {}
        {\begin{itemize}
\item Participation à la collecte et à la distribution de denrées alimentaires.
\item Aide à l'organisation des stocks et à la préparation des campagnes hivernales.
\item Sensibilisation des équipes bénévoles aux bonnes pratiques logistiques.
\end{itemize}}
    
\end{document}